% =====================================================================
%  DROP-IN: new \paragraph for the "Open problems and prospects"
%  subsection (sec:spacetime-open-problems) of
%  spacetime_from_non-computability.tex.
%
%  Drafted 2026-06-13. Adds the two future-work suggestions discussed:
%    (1) Feynman loops as closed computational paths (CTC-like), which
%        would give a SECOND route to the parallelisable-sphere constraint
%        and unify the spacetime arc (CTCs -> dimensions) with the
%        interaction arc (entanglement classes -> gauge group).
%    (2) "Parallelisability determines renormalizability" -- a question
%        posed in the division-algebra community's native language,
%        explicitly offered for them to take up.
%
%  Both are flagged as NOT-yet-well-formed open questions -- suggestions
%  with a possible path, not posed problems. This matches the subsection's
%  existing voice (each prospect with its debt named) and the cautious
%  posture of the QED-scope paragraph in the circuits section, which this
%  complements: that paragraph defers the path-integral construction to
%  Paper 2; this one notes what that construction might, in turn, illuminate.
%
%  PLACEMENT: append after the existing prospect paragraphs (e.g. after
%  "Expansion and acceleration." and before, or as part of, the closing
%  \medskip summary), inside sec:spacetime-open-problems. Adjust the
%  \medskip summary if needed so it still reads as the section's close.
%
%  Cross-refs used (verify): sec:three-dimensions (the CTC ->
%  parallelisable-sphere argument), sec:entanglement-classes (the
%  entanglement-class -> gauge-group route), sec:interactions (the QED-scope
%  paragraph). Citations: adams1960hopf already present; no new keys needed.
% =====================================================================

\paragraph{Loops, supertasks, and a second route to parallelisability.}
The argument of \S\ref{sec:three-dimensions} derives the parallelisable-sphere constraint from
closed timelike curves: a computation whose output is its own input requires information to be
transported consistently around a loop, and consistent transport selects the parallelisable spheres
$S^{1}$, $S^{3}$, $S^{7}$ \autocite{adams1960hopf}. There is a suggestion, which we record as a
prospect rather than a result, that the same constraint may govern the loops of perturbative quantum
field theory. The external legs of a Feynman diagram are on-shell classical boundary data; the
quantum content, and the divergences, enter with the internal loops, where the loop momentum is
integrated over all values and the action accumulated inside the loop is in this sense unbounded ---
a supertask, an infinite computation closed into a finite process. If such a loop is, in the
framework's terms, a closed computational path of the same kind as a closed timelike curve --- a
cycle carrying information not derivable from the data outside it --- then the parallelisability
constraint derived for closed timelike curves would apply to it, and the admissible interactions
would be constrained by the same $S^{1}$, $S^{3}$, $S^{7}$ structure that fixes the gauge group in
\S\ref{sec:entanglement-classes}. This would be a second and independent route to the Standard Model
content: not only ``closed paths in spacetime select the parallelisable spheres'' but ``closed paths
in interactions select them too,'' unifying the spacetime arc and the interaction arc under one
principle. The decisive step, which we have not taken, is the identification of a field-theory loop
with a closed computational path in the precise sense the framework requires --- a loop integral is a
closed cycle in momentum space, a closed timelike curve a closed cycle in spacetime, and whether
these are the same kind of self-consistent closed structure is exactly what would have to be
established. It is offered as a direction, not a claim.

\paragraph{Does parallelisability determine renormalisability?}
A sharper and more self-contained question sits alongside the previous one, and we state it in the
hope that it will be taken up by those equipped for it. Renormalisability is the property that the
divergences of a quantum field theory can be consistently absorbed, leaving finite predictions; some
theories have it and some do not, and the standard criterion for which is technical (power-counting
of the coupling's dimension). The framework suggests a more structural reading. Parallelisable
spheres are precisely the spheres on which a global, consistent transport exists where it generically
does not; the non-parallelisable spheres are where the obstruction lives. If the divergence of a loop
is, in the framework's language, an obstruction to closing a computational path consistently, then
parallelisability would be exactly the structure that distinguishes the removable obstruction (a
renormalisable theory) from the irremovable one. It is at least a suggestive coincidence that the
renormalisable interactions of nature are the gauge theories built on $U(1)$, $SU(2)$ and $SU(3)$ ---
the groups the framework ties to the three parallelisable spheres --- and that gravity, which resists
renormalisation, is the sector \S\ref{sec:three-dimensions} ties to the geometry rather than the
algebra. Whether this is structural or accidental we do not know. The question --- whether
parallelisability, in the Hurwitz--Adams sense, is what determines renormalisability --- is posed in
the native language of the division-algebra approach to the Standard Model, and is one we think that
community is best placed to answer. We are careful to claim nothing more than that the question is
well-motivated: this is emphatically not a method for computing renormalised quantities, whose finite
remainders carry the precision agreement with experiment, but a conjecture about which theories admit
such a computation at all. It is, like the previous prospect, a direction with a possible path to an
answer rather than an answer, and it is left open deliberately.
